% !PRINTED 256
\weberpage{256}
%<<CONT
with the sole exception of the case $n=1$, $p=2$. First we have in $G$ the substitutions
\[
\begin{pmatrix} 0, & -1 \\ 1, & 0 \end{pmatrix}\begin{pmatrix} 0, & 1 \\ -1, & \xi \end{pmatrix}\begin{pmatrix} 0, & 1 \\ -1, & 0 \end{pmatrix} = \begin{pmatrix} \xi, & 1 \\ -1, & 0 \end{pmatrix},
\]
hence also
\[
\begin{pmatrix} 0, & 1 \\ -1, & \xi \end{pmatrix}\begin{pmatrix} \xi, & 1 \\ -1, & 0 \end{pmatrix} = \begin{pmatrix} 1, & 0 \\ 2\xi, & 1 \end{pmatrix},
\]
and consequently
\[
\begin{pmatrix} 1, & 0 \\ 2\xi, & 1 \end{pmatrix}\begin{pmatrix} 0, & 1 \\ -1, & \xi \end{pmatrix} = \begin{pmatrix} 0, & 1 \\ -1, & 3\xi \end{pmatrix},
\]
\[
\begin{pmatrix} 1, & 0 \\ 2\xi, & 1 \end{pmatrix}\begin{pmatrix} 0, & 1 \\ -1, & 3\xi \end{pmatrix} = \begin{pmatrix} 0, & 1 \\ -1, & 5\xi \end{pmatrix} \ \text{and so on,}
\]
\[
\begin{pmatrix} 0, & 1 \\ -1, & m\xi \end{pmatrix}
\]
for every odd integer $m$; if therefore $p$ is odd, we can take $m=p$ and obtain $A_0$. But if $p=2$, then, denoting by $\varrho$ a number in $\mathfrak{C}$ still to be determined, we have the following substitutions in $G$:
\[
\begin{pmatrix} 0, & \varrho \\ -\varrho^{-1}, & 0 \end{pmatrix}\begin{pmatrix} 0, & 1 \\ -1, & \xi \end{pmatrix}\begin{pmatrix} 0, & -\varrho \\ \varrho^{-1}, & 0 \end{pmatrix} = \begin{pmatrix} \xi, & \varrho^2 \\ -\varrho^{-2}, & 0 \end{pmatrix}.
\]
Since here every number is a square, one can also replace $\varrho^2$ by $\varrho$, and thus finds that in $G$ the substitution
\[
\begin{pmatrix} \xi, & \varrho \\ -\varrho^{-1}, & 0 \end{pmatrix}
\]
for every nonzero $\varrho$, and consequently also its inverse
\[
\begin{pmatrix} 0, & -\varrho \\ \varrho^{-1}, & \xi \end{pmatrix}
\]
must occur. If, then, $\varrho,\sigma$ denote two numbers different from zero, we have in $G$ also
\weq{9}{\begin{pmatrix} 0, & \varrho^{-1}\sigma^{-1} \\ -\varrho\sigma, & \xi \end{pmatrix}\begin{pmatrix} \xi, & \sigma^{-1} \\ -\sigma, & 0 \end{pmatrix}\begin{pmatrix} 0, & \varrho \\ -\varrho^{-1}, & \xi \end{pmatrix} = \begin{pmatrix} 0, & 1 \\ -1, & \xi\varrho(1+\sigma+\sigma\varrho) \end{pmatrix}.}
